\documentclass[10pt,a4paper]{amsart}
\usepackage{amsmath,amssymb,amsthm,graphicx}
\usepackage{verbatim}
\newenvironment{reference}{\comment}{\endcomment}
%\newenvironment{reference}{}{}
\newenvironment{slogan}{\comment}{\endcomment}
\newenvironment{history}{\comment}{\endcomment}

% For commutative diagrams we use Xy-pic
\usepackage[all]{xy}

% We use 2cell for 2-commutative diagrams.
\xyoption{2cell}
% Theorem environments.
%
\theoremstyle{plain}
\newtheorem{theorem}[subsection]{Theorem}
\newtheorem{proposition}[subsection]{Proposition}
\newtheorem{lemma}[subsection]{Lemma}

\theoremstyle{definition}
\newtheorem{definition}[subsection]{Definition}
\newtheorem{example}[subsection]{Example}
\newtheorem{exercise}[subsection]{Exercise}
\newtheorem{situation}[subsection]{Situation}

\theoremstyle{remark}
\newtheorem{remark}[subsection]{Remark}
\newtheorem{remarks}[subsection]{Remarks}

\numberwithin{equation}{subsection}

% Macros
%
\def\lim{\mathop{\mathrm{lim}}\nolimits}
\def\colim{\mathop{\mathrm{colim}}\nolimits}
\def\Spec{\mathop{\mathrm{Spec}}}
\def\Hom{\mathop{\mathrm{Hom}}\nolimits}
\def\Ext{\mathop{\mathrm{Ext}}\nolimits}
\def\SheafHom{\mathop{\mathcal{H}\!\mathit{om}}\nolimits}
\def\SheafExt{\mathop{\mathcal{E}\!\mathit{xt}}\nolimits}
\def\Sch{\mathit{Sch}}
\def\Mor{\mathop{\mathrm{Mor}}\nolimits}
\def\Ob{\mathop{\mathrm{Ob}}\nolimits}
\def\Sh{\mathop{\mathit{Sh}}\nolimits}
\def\NL{\mathop{N\!L}\nolimits}
\def\CH{\mathop{\mathrm{CH}}\nolimits}
\def\proetale{{pro\text{-}\acute{e}tale}}
\def\etale{{\acute{e}tale}}
\def\QCoh{\mathit{QCoh}}
\def\Ker{\mathop{\mathrm{Ker}}}
\def\Im{\mathop{\mathrm{Im}}}
\def\Coker{\mathop{\mathrm{Coker}}}
\def\Coim{\mathop{\mathrm{Coim}}}

% Boxtimes
%
\DeclareMathSymbol{\boxtimes}{\mathbin}{AMSa}{"02}

%
% Macros for moduli stacks/spaces
%
\def\QCohstack{\mathcal{QC}\!\mathit{oh}}
\def\Cohstack{\mathcal{C}\!\mathit{oh}}
\def\Spacesstack{\mathcal{S}\!\mathit{paces}}
\def\Quotfunctor{\mathrm{Quot}}
\def\Hilbfunctor{\mathrm{Hilb}}
\def\Curvesstack{\mathcal{C}\!\mathit{urves}}
\def\Polarizedstack{\mathcal{P}\!\mathit{olarized}}
\def\Complexesstack{\mathcal{C}\!\mathit{omplexes}}
% \Pic is the operator that assigns to X its picard group, usage \Pic(X)
% \Picardstack_{X/B} denotes the Picard stack of X over B
% \Picardfunctor_{X/B} denotes the Picard functor of X over B
\def\Pic{\mathop{\mathrm{Pic}}\nolimits}
\def\Picardstack{\mathcal{P}\!\mathit{ic}}
\def\Picardfunctor{\mathrm{Pic}}
\def\Deformationcategory{\mathcal{D}\!\mathit{ef}}
\title{When the pushforward of abelian sheaves reflects surjections, and the resulting faithfulness}
\author{The Stacks Project Authors}
\date{stacks-project:sites-modules.tex:local-snapshot-20261012; GNU FDL 1.2 or later}
\begin{document}
\maketitle
\noindent\textit{Editorial note.} All delivered mathematical text is the author's own, copied verbatim from the pinned Stacks Project chapter file sites-modules.tex; no statement or proof step was rewritten, reordered or removed.\par\medskip
\noindent\textit{Editorial note.} The counted span cites no labelled external item of this chapter (the exactness of f\textasciicircum{}\{-1\} is discussed in the sites chapter), so no navigation stubs were needed.\par\medskip
\noindent\textit{Allowed standard prerequisites (editorial).}\begin{itemize}
\item Exactness of the pullback functor f\textasciicircum{}\{-1\} on abelian sheaves. Stacks Project, Sites, section on morphisms of topoi (exactness of f\textasciicircum{}\{-1\}). Justifies transporting surjectivity from f\_*a to f\textasciicircum{}\{-1\}f\_*a.
\end{itemize}\medskip
\begin{lemma}
\label{lemma-reflect-surjections}
Let $f : \Sh(\mathcal{C}) \to \Sh(\mathcal{D})$ be
a morphism of topoi. The following are equivalent:
\begin{enumerate}
\item $f^{-1}f_*\mathcal{F} \to \mathcal{F}$ is surjective for
all $\mathcal{F}$ in $\textit{Ab}(\mathcal{C})$, and
\item $f_* : \textit{Ab}(\mathcal{C}) \to \textit{Ab}(\mathcal{D})$
reflects surjections.
\end{enumerate}
In this case the functor
$f_* : \textit{Ab}(\mathcal{C}) \to \textit{Ab}(\mathcal{D})$
is faithful.
\end{lemma}

\begin{proof}
Assume (1). Suppose that $a : \mathcal{F} \to \mathcal{F}'$
is a map of abelian sheaves on $\mathcal{C}$ such that $f_*a$ is surjective.
As $f^{-1}$ is exact this implies that
$f^{-1}f_*a : f^{-1}f_*\mathcal{F} \to f^{-1}f_*\mathcal{F}'$
is surjective. Combined with (1) this implies that $a$ is surjective.
This means that (2) holds.

\medskip\noindent
Assume (2). Let $\mathcal{F}$ be an abelian sheaf on $\mathcal{C}$.
We have to show that the map $f^{-1}f_*\mathcal{F} \to \mathcal{F}$ is
surjective. By (2) it suffices to show that
$f_*f^{-1}f_*\mathcal{F} \to f_*\mathcal{F}$ is surjective.
And this is true because there is a canonical map
$f_*\mathcal{F} \to f_*f^{-1}f_*\mathcal{F}$ which is a one-sided inverse.

\medskip\noindent
We omit the proof of the final assertion.
\end{proof}

\end{document}
